\chapter{Analysis \& Model Architecture — Comprehensive Reference}
\label{ch:analysis_model_architecture}

% =============================================================================
%  Chapter 63: Full architectural documentation of the Analysis and Model
%  modules, including dependency maps, extension points, usage tutorials,
%  and semantic critique.
% =============================================================================


% ═══════════════════════════════════════════════════════════════════════════
\section{Overview}
\label{sec:arch:overview}

The \texttt{fall\_n} framework organises finite-element analysis into two
principal modules:

\begin{description}[style=nextline]
  \item[\texttt{src/model/}]
    Describes \emph{what} is being analysed: the finite element model, its
    degrees of freedom, boundary conditions, material points, and state
    snapshots.

  \item[\texttt{src/analysis/}]
    Describes \emph{how} the model is solved: linear, nonlinear, and dynamic
    solvers; incremental load control; step-level observation; and multi-phase
    orchestration.
\end{description}

Both modules are policy-based (templates) at the inner core and type-erased
(callbacks, \texttt{std::function}) at the system boundary.  This gives
zero-cost inner loops and maximum composability at the API surface.

The next sections dissect each module by:
\begin{itemize}[nosep]
  \item Header inventory and dependency graph,
  \item Class-level architecture diagrams,
  \item Extension points with concrete tutorials,
  \item Semantic analysis and naming conventions.
\end{itemize}


% ═══════════════════════════════════════════════════════════════════════════
% ═══════════════════════════════════════════════════════════════════════════
\section{Analysis Module Architecture}
\label{sec:arch:analysis}


% ─────────────────────────────────────────────────────────────────────────
\subsection{Header Inventory}
\label{sec:arch:analysis:inventory}

The analysis module comprises 13 header files, all flat in
\texttt{src/analysis/} (no subdirectories):

\begin{center}
\small
\begin{tabular}{@{}llp{7.5cm}@{}}
\toprule
\textbf{Header} & \textbf{Lines} & \textbf{Role} \\
\midrule
\texttt{AnalysisObserver.hh}     & $\sim$320 & \texttt{StepEvent}, \texttt{ObserverBase}, \texttt{DynamicObserverList}, \texttt{CompositeObserver}, \texttt{ObserverCallback} \\
\texttt{IncrementalControl.hh}   & $\sim$170 & \texttt{IncrementalControlPolicy} concept $+$ \texttt{LoadControl}, \texttt{DisplacementControl}, \texttt{CustomControl} \\
\texttt{StepDirector.hh}         & $\sim$180 & \texttt{StepVerdict} enum, \texttt{StepDirector<ModelT>} type alias, 5~factories \\
\texttt{SteppableSolver.hh}      & $\sim$70  & \texttt{SteppableSolver} concept $+$ \texttt{AnalysisState} struct \\
\texttt{Damping.hh}              & $\sim$100 & \texttt{DampingAssembler} functor $+$ Rayleigh factories \\
\texttt{LinearAnalysis.hh}       & $\sim$95  & PETSc KSP wrapper (single solve) \\
\texttt{NLAnalysis.hh}           & $\sim$870 & PETSc SNES Newton--Raphson solver with automatic bisection \\
\texttt{DynamicAnalysis.hh}      & $\sim$1000& PETSc TS transient solver (generalized-$\alpha$, Newmark) \\
\texttt{AnalysisDirector.hh}     & $\sim$280 & Multi-phase orchestrator with state threading \\
\texttt{Observers.hh}            & $\sim$800 & 7~concrete observers (console, VTK, node/element/fiber recorders, energy, max-response) \\
\texttt{DamageCriterion.hh}      & $\sim$320 & Damage-based stopping criterion (observer) \\
\texttt{FiberHysteresisRecorder.hh} & $\sim$330 & Fiber-level hysteresis recorder (observer) \\
\texttt{Analysis.hh}             & $\sim$50  & Module aggregator (includes all of the above in dependency order) \\
\bottomrule
\end{tabular}
\end{center}


% ─────────────────────────────────────────────────────────────────────────
\subsection{Dependency Graph}
\label{sec:arch:analysis:deps}

The include relationships form a layered DAG:

\begin{center}
\begin{tikzpicture}[
    box/.style={draw, rounded corners=2pt, minimum width=3.8cm,
                minimum height=0.8cm, font=\footnotesize\ttfamily,
                align=center},
    concept/.style={box, fill=blue!10},
    solver/.style={box, fill=green!10},
    orch/.style={box, fill=orange!10},
    obs/.style={box, fill=yellow!10},
    facade/.style={box, fill=gray!10},
    arr/.style={-{Stealth[length=4pt]}, thin, gray!70}
  ]

  % Layer 0: Model
  \node[box, fill=red!8] (model) at (0, 6)
    {Model.hh\\(src/model)};

  % Layer 1: Concepts / value types
  \node[concept] (observer) at (-5.5, 4.5)   {AnalysisObserver.hh};
  \node[concept] (incctrl)  at (-1.8, 4.5)   {IncrementalControl.hh};
  \node[concept] (stepdir)  at ( 1.8, 4.5)   {StepDirector.hh};
  \node[concept] (steppable)at ( 5.5, 4.5)   {SteppableSolver.hh};

  % Layer 1.5: Damping
  \node[concept] (damping)  at (-5.5, 3.3)   {Damping.hh};

  % Layer 2: Solvers
  \node[solver] (linear) at (-4, 1.8)   {LinearAnalysis.hh};
  \node[solver] (nl)     at ( 0, 1.8)   {NLAnalysis.hh};
  \node[solver] (dyn)    at ( 4, 1.8)   {DynamicAnalysis.hh};

  % Layer 3: Orchestration
  \node[orch] (director) at (0, 0.3) {AnalysisDirector.hh};

  % Layer 4: Observers
  \node[obs] (observers) at (-4, -1.0) {Observers.hh};
  \node[obs] (damage)    at ( 0, -1.0) {DamageCriterion.hh};
  \node[obs] (fiber)     at ( 4, -1.0) {FiberHysteresisRecorder.hh};

  % Layer 5: Facade
  \node[facade] (analysis) at (0, -2.5) {Analysis.hh};

  % Edges
  \draw[arr] (model) -- (observer);
  \draw[arr] (model) -- (incctrl);
  \draw[arr] (observer) -- (stepdir);
  \draw[arr] (stepdir) -- (steppable);

  \draw[arr] (model) -- (linear);
  \draw[arr] (model) -- (nl);
  \draw[arr] (model) -- (dyn);
  \draw[arr] (incctrl) -- (nl);
  \draw[arr] (stepdir) -- (nl);
  \draw[arr] (damping) -- (dyn);
  \draw[arr] (observer) -- (dyn);
  \draw[arr] (stepdir) -- (dyn);
  \draw[arr] (steppable) -- (dyn);

  \draw[arr] (steppable) -- (director);

  \draw[arr] (observer) -- (observers);
  \draw[arr] (observer) -- (damage);
  \draw[arr] (observer) -- (fiber);

  \draw[arr] (linear) -- (analysis);
  \draw[arr] (nl) -- (analysis);
  \draw[arr] (dyn) -- (analysis);
  \draw[arr] (director) -- (analysis);
  \draw[arr] (observers) -- (analysis);
  \draw[arr] (damage) -- (analysis);
  \draw[arr] (fiber) -- (analysis);

\end{tikzpicture}
\end{center}

\texttt{Analysis.hh} is a pure aggregator: it includes everything in
dependency order so that a single \texttt{\#include "Analysis.hh"} brings
the entire module into scope.


% ─────────────────────────────────────────────────────────────────────────
\subsection{Class-Level Architecture}
\label{sec:arch:analysis:classes}

The analysis module is organised around four design pillars:

\begin{enumerate}
  \item \textbf{Concepts as contracts.}
        \texttt{IncrementalControlPolicy}, \texttt{SteppableSolver},
        and \texttt{StepDirector<ModelT>} are not base classes ---
        they are C++23 concepts or type aliases for \texttt{std::function}.
        Solvers satisfy these contracts structurally; no inheritance is
        required.

  \item \textbf{Solver engines as monolithic templates.}
        \texttt{LinearAnalysis}, \texttt{NonlinearAnalysis}, and
        \texttt{DynamicAnalysis} are class templates parametrised by
        \texttt{MaterialPolicy}, \texttt{KinematicPolicy}, \texttt{ndofs},
        and \texttt{ElemPolicy}.  Each wraps a PETSc solver object (KSP,
        SNES, TS) with RAII ownership.

  \item \textbf{Observer pattern for non-intrusive recording.}
        Observers decouple data recording from the solver loop.  Two
        flavours exist: compile-time (\texttt{CompositeObserver}) and
        runtime (\texttt{DynamicObserverList}).  Both are integrated via
        a single \texttt{ObserverCallback} trampoline.

  \item \textbf{Orchestration at the top.}
        \texttt{AnalysisDirector} composes solver phases using
        type-erased callbacks (\texttt{PhaseCallback}).  State flows
        between phases through \texttt{AnalysisState}.
\end{enumerate}


% ─────────────────────────────────────────────────────────────────────────
\subsection{Solver Engine Anatomy}
\label{sec:arch:analysis:solvers}

All three solver engines share a common internal structure:

\begin{center}
\small
\begin{tabular}{@{}lp{3.5cm}p{3.5cm}p{3.5cm}@{}}
\toprule
\textbf{Aspect} & \textbf{LinearAnalysis} & \textbf{NLAnalysis} & \textbf{DynamicAnalysis} \\
\midrule
PETSc solver   & KSP              & SNES              & TS \\
Equation       & $Ku = f$         & $R(u) = f_{\text{int}}(u) - f_{\text{ext}} = 0$ & $M\ddot{u} + C\dot{u} + f_{\text{int}}(u) = f_{\text{ext}}(t)$ \\
Callbacks      & none             & \texttt{FormResidual}, \texttt{FormJacobian} & \texttt{I2Function}, \texttt{I2Jacobian}, \texttt{Monitor} \\
Stepping       & single solve     & incremental $+$ bisection & time integration \\
State commit   & post-solve       & per converged step & per converged step \\
Observer       & none             & \texttt{step\_callback\_} (legacy) & \texttt{ObserverCallback} $+$ \texttt{MonitorCallback} \\
Single-step    & no               & \texttt{step()}/\texttt{step\_to()}/\texttt{step\_n()} & same \\
SteppableSolver& no               & yes                & yes \\
\bottomrule
\end{tabular}
\end{center}

\paragraph{NonlinearAnalysis internals.}
The file is organised as follows:
\begin{enumerate}[nosep]
  \item \emph{Private members:} model pointer, PETSc objects (SNES, Vec, Mat),
        timer, callbacks, single-step state, context struct.
  \item \emph{Static PETSc callbacks:} \texttt{FormResidual}, \texttt{FormJacobian}
        --- these are \texttt{static} member functions matching PETSc's C API.
        The \texttt{Context} struct provides the \texttt{void*} bridge.
  \item \emph{State management:} \texttt{commit\_state()}, \texttt{revert\_state()}.
  \item \emph{Public query API:} \texttt{converged\_reason()},
        \texttt{num\_iterations()}, \texttt{timer()}.
  \item \emph{Setup.}
  \item \emph{Solve API:} \texttt{solve()}, \texttt{solve\_incremental()}.
  \item \emph{Bisection engine:} \texttt{advance\_to\_p\_()},
        \texttt{advance\_step\_impl\_()}.
  \item \emph{Single-step API:} \texttt{begin\_incremental()}, \texttt{step()},
        \texttt{step\_to()}, \texttt{step\_n()}.
  \item \emph{Accessors:} \texttt{current\_time()}, \texttt{current\_step()}.
  \item \emph{State transfer:} \texttt{get\_analysis\_state()}.
  \item \emph{Constructor/destructor.}
\end{enumerate}

\paragraph{DynamicAnalysis internals.}
\begin{enumerate}[nosep]
  \item \emph{Private members:} model pointer, PETSc objects (TS, Mat $\times$4, Vec $\times$4),
        configuration callbacks, observer pipelines, ground-motion data.
  \item \emph{Static PETSc callbacks:} \texttt{I2Function}, \texttt{I2Jacobian},
        \texttt{Monitor}.
  \item \emph{Private helpers:} \texttt{post\_step\_()},
        \texttt{assemble\_initial\_stiffness()},
        \texttt{setup\_ground\_motion\_influence()}.
  \item \emph{Configuration API:} density, damping, forces, BCs, observers.
  \item \emph{Setup.}
  \item \emph{Initial conditions.}
  \item \emph{Single-step API:} \texttt{step()}, \texttt{step\_to()}, \texttt{step\_n()}.
  \item \emph{Runtime reconfiguration:} \texttt{set\_time\_step()}, \texttt{set\_director()}.
  \item \emph{Batch solve:} \texttt{solve(t\_final, dt)}, \texttt{solve(t0, t1, dt)}.
  \item \emph{Post-solve accessors:} displacement, velocity, matrices, \texttt{get\_analysis\_state()}.
  \item \emph{Constructor/destructor.}
\end{enumerate}


% ─────────────────────────────────────────────────────────────────────────
\subsection{Incremental Control Policies}
\label{sec:arch:analysis:control}

The incremental solver in \texttt{NonlinearAnalysis} accepts any control
strategy satisfying:

\begin{verbatim}
template <typename S, typename ModelT>
concept IncrementalControlPolicy = requires(
    S& scheme, double p, Vec f_full, Vec f_ext, ModelT* model)
{
    scheme.apply(p, f_full, f_ext, model);
};
\end{verbatim}

\noindent
The \texttt{apply()} method is \textbf{absolute}: for a given $p$, it fully
determines the system state.  This idempotency is critical for bisection
rollback --- no save/restore logic is needed in the scheme itself.

Three built-in policies are provided:

\begin{center}
\begin{tabular}{@{}lp{8cm}@{}}
\toprule
\textbf{Policy} & \textbf{Semantics} \\
\midrule
\texttt{LoadControl}         & $f_{\text{ext}} = p \cdot f_{\text{full}}$ (proportional scaling, default) \\
\texttt{DisplacementControl} & Prescribed DOF $u_D = p \cdot u_{\text{target}}$; optional $f_{\text{ext}} = p \cdot \alpha \cdot f_{\text{full}}$ \\
\texttt{CustomControl<F>}    & User-supplied $\lambda(p)$ via \texttt{make\_control(fn)} \\
\bottomrule
\end{tabular}
\end{center}

\noindent
\textbf{Planned:} \texttt{ArcLengthControl} (Riks method) requires augmenting
the SNES system to $(u, \lambda)$ with a spherical constraint.


% ─────────────────────────────────────────────────────────────────────────
\subsection{Observer Architecture}
\label{sec:arch:analysis:observers}

\begin{center}
\begin{tikzpicture}[
    box/.style={draw, rounded corners=2pt, minimum width=3.5cm,
                minimum height=0.7cm, font=\small\ttfamily},
    arr/.style={-{Stealth[length=4pt]}, thick}
  ]
  \node[box, fill=blue!10] (base) at (0, 3)
    {ObserverBase<ModelT>};
  \node[box, fill=green!10] (dyn) at (-3.5, 1.5)
    {DynamicObserverList};
  \node[box, fill=green!10] (comp) at (3.5, 1.5)
    {CompositeObserver};
  \node[box, fill=orange!10] (cb) at (0, 0)
    {ObserverCallback<ModelT>};

  \draw[arr] (base) -- (dyn) node[midway, left, font=\scriptsize] {owns via unique\_ptr};
  \draw[arr] (comp.south) -- (cb) node[midway, right, font=\scriptsize] {make\_observer\_callback};
  \draw[arr] (dyn.south) -- (cb)  node[midway, left, font=\scriptsize]  {make\_observer\_callback};
\end{tikzpicture}
\end{center}

\begin{itemize}[nosep]
  \item \textbf{Static path} (\texttt{CompositeObserver}): observers stored in
        a \texttt{std::tuple}, fold-expanded at compile time.  Zero virtual
        dispatch.  Use when the observer set is fixed at compile time.

  \item \textbf{Dynamic path} (\texttt{DynamicObserverList}): observers behind
        \texttt{unique\_ptr<ObserverBase>}, virtual dispatch per step.
        Negligible cost compared to constitutive evaluation.  Use when
        observers are configured at runtime.

  \item \textbf{Integration}: \texttt{DynamicAnalysis::set\_observer()} stores
        an \texttt{ObserverCallback} --- a struct of three
        \texttt{std::function} fields (\texttt{on\_start}, \texttt{on\_step},
        \texttt{on\_end}).  The \texttt{make\_observer\_callback<ModelT>(obs)}
        factory bridges any observer-like object to this interface.
\end{itemize}

\paragraph{Concrete observers.}

\begin{center}
\small
\begin{tabular}{@{}lp{8cm}@{}}
\toprule
\textbf{Observer} & \textbf{Purpose} \\
\midrule
\texttt{ConsoleProgressObserver} & Print step/time to stdout every $N$ steps \\
\texttt{VTKSnapshotObserver}     & Write VTU/VTM files at intervals \\
\texttt{NodeRecorder}            & Record displacement/velocity at selected nodes \\
\texttt{ElementRecorder}         & Record element-level fields \\
\texttt{FiberRecorder}           & Record fiber-level $\sigma$--$\varepsilon$ history \\
\texttt{EnergyRecorder}          & Track kinetic, strain, and external work \\
\texttt{MaxResponseTracker}      & Track peak displacement/acceleration \\
\texttt{DamageCriterion}         & Evaluate damage-based stopping criteria \\
\texttt{FiberHysteresisRecorder} & Record fiber-level hysteresis loops \\
\bottomrule
\end{tabular}
\end{center}


% ─────────────────────────────────────────────────────────────────────────
\subsection{StepDirector and SteppableSolver}
\label{sec:arch:analysis:stepping}

Both \texttt{NonlinearAnalysis} and \texttt{DynamicAnalysis} satisfy the
\texttt{SteppableSolver} concept (Chapter~\ref{ch:steppable_solver_concept}).
The ``time axis'' has engine-specific semantics:

\begin{center}
\begin{tabular}{@{}lll@{}}
\toprule
\textbf{Engine} & \textbf{Time axis} & \textbf{Range} \\
\midrule
\texttt{NonlinearAnalysis} & Control parameter $p$ & $[0, 1]$ \\
\texttt{DynamicAnalysis}   & Physical time $t$     & $[0, T]$ \\
\bottomrule
\end{tabular}
\end{center}

The \texttt{StepDirector<ModelT>} is a
\texttt{std::function<StepVerdict(const StepEvent\&, const ModelT\&)>}
evaluated after each converged step.  Five factories cover common patterns:

\begin{enumerate}[nosep]
  \item \texttt{pause\_at\_times} --- pause at listed breakpoints.
  \item \texttt{pause\_every\_n} --- pause every $N$ steps.
  \item \texttt{pause\_on(pred)} --- pause when a predicate fires.
  \item \texttt{stop\_on(pred)} --- terminate on a predicate.
  \item \texttt{compose} --- combine directors (most restrictive wins).
\end{enumerate}


% ─────────────────────────────────────────────────────────────────────────
\subsection{AnalysisDirector (Orchestration)}
\label{sec:arch:analysis:director}

Multi-phase workflows are coordinated by \texttt{AnalysisDirector}
(Chapter~\ref{ch:analysis_director}).  Key design choices:

\begin{itemize}[nosep]
  \item Director is \emph{not} templated --- phases are type-erased
        \texttt{PhaseCallback = std::function<AnalysisState(const AnalysisState\&)>}.
  \item State flows automatically: the output of phase $n$ becomes the input
        of phase $n{+}1$.
  \item Optional \texttt{PhaseGateCallback} for inter-phase decisions.
  \item Returns a \texttt{DirectorReport} with per-phase timing and status.
\end{itemize}


% ═══════════════════════════════════════════════════════════════════════════
% ═══════════════════════════════════════════════════════════════════════════
\section{Model Module Architecture}
\label{sec:arch:model}


% ─────────────────────────────────────────────────────────────────────────
\subsection{Header Inventory}
\label{sec:arch:model:inventory}

The model module comprises 12 header files in \texttt{src/model/}:

\begin{center}
\small
\begin{tabular}{@{}llp{7.5cm}@{}}
\toprule
\textbf{Header} & \textbf{Lines} & \textbf{Role} \\
\midrule
\texttt{DoFStorage.hh}        & $\sim$400 & 3~DOF storage policies $+$ concept; default used by \texttt{Domain} \\
\texttt{DoF.hh}               & $\sim$130 & Legacy DOF handler (superseded by \texttt{DoFStorage}; retained for transition) \\
\texttt{MaterialPoint.hh}     & $\sim$130 & Lightweight spatial site carrying material reference \\
\texttt{ModelState.hh}         & $\sim$140 & Serialisable state snapshot (displacement, velocity, time) \\
\texttt{NodeSelector.hh}       & $\sim$330 & Concept-based composable node selection (by id, label, box, functor) \\
\texttt{Model.hh}              & $\sim$550 & Core FEA model: elements, PETSc section, constraints, forces \\
\texttt{ModelBuilder.hh}       & $\sim$80  & \emph{Dormant} --- original domain$\leftrightarrow$model linker (all code commented) \\
\texttt{BoundaryCondition.hh}  & $\sim$690 & Time-dependent boundary conditions, IC, ground motion \\
\texttt{LoadUtilities.hh}      & $\sim$50  & Helper to apply nodal loads \\
\texttt{BuildingDomainBuilder.hh} & $\sim$350 & Generates multi-story regular building meshes \\
\texttt{StructuralModelBuilder.hh} & $\sim$150 & High-level builder for beam/shell structural models \\
\texttt{GroundMotionRecord.hh} & $\sim$280 & Parser for seismic acceleration records (CSV/column) \\
\bottomrule
\end{tabular}
\end{center}


% ─────────────────────────────────────────────────────────────────────────
\subsection{The \texttt{Model<>} Template}
\label{sec:arch:model:model}

\texttt{Model} is the central aggregate: it owns the PETSc data structures
and connects the domain, elements, materials, and boundary conditions.

\begin{verbatim}
template <
    typename MaterialPolicy,
    typename KinematicPolicy = continuum::SmallStrain,
    std::size_t ndofs = MaterialPolicy::dim,
    typename ElemPolicy = SingleElementPolicy<
        ContinuumElement<MaterialPolicy, ndofs, KinematicPolicy>>
>
class Model;
\end{verbatim}

\noindent
Template parameters:

\begin{description}[style=nextline, leftmargin=2cm]
  \item[\texttt{MaterialPolicy}]
    Defines the constitutive dimension (\texttt{dim}), strain/stress types.
    Examples: \texttt{ThreeDimensionalMaterial}, \texttt{PlaneMaterial},
    \texttt{UniaxialMaterial}.

  \item[\texttt{KinematicPolicy}]
    Selects the kinematic formulation.
    Options: \texttt{SmallStrain}, \texttt{TotalLagrangian},
    \texttt{UpdatedLagrangian}, \texttt{Corotational}.

  \item[\texttt{ndofs}]
    Degrees of freedom per node (defaults to \texttt{dim} for solid
    elements; 6~for structural beam/shell).

  \item[\texttt{ElemPolicy}]
    Storage policy for elements.  \texttt{SingleElementPolicy} stores a
    homogeneous \texttt{std::vector<ElementT>};
    \texttt{MultiElementPolicy} stores \texttt{std::vector<FEM\_Element>}
    (type-erased, heterogeneous).
\end{description}

\paragraph{Responsibilities.}
\begin{enumerate}[nosep]
  \item \textbf{PETSc section management}: creates and configures the
        \texttt{PetscSection} on the DMPlex, assigning DOF counts per
        mesh point.
  \item \textbf{Element storage}: owns the element container; provides
        \texttt{elements()} accessor for iteration.
  \item \textbf{Global vectors}: \texttt{force\_vector()},
        \texttt{imposed\_solution()}, \texttt{state\_vector()} ---
        PETSc local vectors managed via RAII.
  \item \textbf{Constraint management}: \texttt{constrain\_dof()},
        \texttt{update\_imposed\_value()} --- manipulates the DM section
        to enforce essential BCs.
  \item \textbf{Stiffness injection}: \texttt{inject\_K(Mat)} --- assembles
        the global tangent stiffness by iterating over elements.
  \item \textbf{Mass assembly}: \texttt{assemble\_mass\_matrix(Mat)} for
        dynamics.
  \item \textbf{State capture}: \texttt{capture\_state()} /
        \texttt{restore\_state()} for checkpoint/restart.
\end{enumerate}


% ─────────────────────────────────────────────────────────────────────────
\subsection{Boundary Condition Framework}
\label{sec:arch:model:bc}

The \texttt{BoundaryConditionSet<dim>} class provides a declarative API
for time-dependent boundary conditions:

\begin{itemize}[nosep]
  \item \textbf{Nodal forces:} point loads with time modulation.
  \item \textbf{Prescribed displacements:} essential BCs with time functions.
  \item \textbf{Initial conditions:} displacement and velocity ICs.
  \item \textbf{Ground motions:} seismic input (direction + acceleration record).
\end{itemize}

\noindent
Integration with \texttt{DynamicAnalysis}:

\begin{verbatim}
BoundaryConditionSet<3> bcs;
bcs.add_force(node_selector, {0, 0, -1000}, linear_ramp);
bcs.add_prescribed(base_nodes, 0, ground_motion_fn);
bcs.add_ground_motion(0, accel_fn);

dyn.set_boundary_conditions(bcs);
dyn.set_initial_conditions(bcs);
\end{verbatim}


% ─────────────────────────────────────────────────────────────────────────
\subsection{NodeSelector: Composable Node Queries}
\label{sec:arch:model:selector}

\texttt{NodeSelector} uses a concept + combinator pattern:

\begin{verbatim}
template <typename S>
concept NodeSelectorConcept = requires(const S& sel, const Node& n) {
    { sel(n) } -> std::convertible_to<bool>;
};
\end{verbatim}

Built-in selectors:
\begin{itemize}[nosep]
  \item \texttt{ByNodeId\{ids\}} --- select by explicit node ids.
  \item \texttt{ByLabel\{label\}} --- select by physical group label.
  \item \texttt{ByBoundingBox\{min, max\}} --- spatial selection.
  \item \texttt{ByPredicate\{fn\}} --- arbitrary predicate.
\end{itemize}

Combinators: \texttt{operator\&\&}, \texttt{operator||}, \texttt{operator!}
produce composed selectors.  All are value types (no heap allocation).


% ─────────────────────────────────────────────────────────────────────────
\subsection{Semantic Analysis}
\label{sec:arch:model:semantics}

\begin{description}[style=nextline]
  \item[\texttt{DoF.hh} vs.\ \texttt{DoFStorage.hh}]
    The legacy \texttt{DoF\_Interface}/\texttt{DoF\_Handler} classes use
    \texttt{shared\_ptr} and a mutable public member style inconsistent
    with the rest of the framework.  \texttt{DoFStorage.hh} replaces them
    with three clean storage policies (\texttt{FixedDoFStorage},
    \texttt{DynamicDoFStorage}, \texttt{PetscDoFStorage}).  The old header
    is retained only for downstream compatibility and should be removed
    when all consumers migrate.

  \item[\texttt{ModelBuilder.hh}]
    All code is commented out.  The original purpose (linking Domain to
    Model) has been absorbed by \texttt{StructuralModelBuilder} and
    \texttt{GmshDomainBuilder}.  The empty shell is kept to avoid breaking
    \texttt{header\_files.hh} includes.

  \item[\texttt{MaterialPoint.hh}]
    Correctly named: it represents a spatial location carrying a material
    reference (a Gauss point with material identity).

  \item[\texttt{ModelState.hh}]
    Clean naming; captures a serialisable snapshot of the model.
\end{description}


% ═══════════════════════════════════════════════════════════════════════════
% ═══════════════════════════════════════════════════════════════════════════
\section{Extension Points}
\label{sec:arch:extension}


% ─────────────────────────────────────────────────────────────────────────
\subsection{Adding a New Solver Engine}
\label{sec:arch:ext:solver}

To add a new analysis engine (e.g., explicit dynamics or eigenvalue analysis):

\begin{enumerate}
  \item \textbf{Create the header} in \texttt{src/analysis/}.  Use the same
        template signature (\texttt{MaterialPolicy}, \texttt{KinematicPolicy},
        \texttt{ndofs}, \texttt{ElemPolicy}) for compatibility.

  \item \textbf{Wrap the PETSc solver} (e.g., \texttt{EPS} for eigenvalue)
        with RAII via \texttt{petsc::OwnedXxx}.

  \item \textbf{Implement the solve API}.  If the engine supports stepping,
        provide \texttt{step()}, \texttt{step\_to()}, \texttt{step\_n()},
        \texttt{current\_time()}, \texttt{current\_step()} to satisfy
        \texttt{SteppableSolver}.

  \item \textbf{Provide \texttt{get\_analysis\_state()}} returning an
        \texttt{AnalysisState} for integration with \texttt{AnalysisDirector}.

  \item \textbf{Register in \texttt{Analysis.hh}} by adding an include in
        the ``Solver engines'' section.
\end{enumerate}

\begin{verbatim}
// Example skeleton:
template <typename MaterialPolicy, ...>
class EigenvalueAnalysis {
    ModelT* model_;
    petsc::OwnedEPS eps_;
public:
    void solve(int num_modes);
    auto eigenvalues() const -> std::vector<double>;
    auto eigenvector(int i) const -> Vec;

    explicit EigenvalueAnalysis(ModelT* model);
};
\end{verbatim}


% ─────────────────────────────────────────────────────────────────────────
\subsection{Adding a New Observer}
\label{sec:arch:ext:observer}

\paragraph{Runtime observer (recommended for most use cases):}

\begin{enumerate}[nosep]
  \item Derive from \texttt{ObserverBase<ModelT>}.
  \item Override \texttt{on\_step()} (required) and optionally
        \texttt{on\_analysis\_start()} / \texttt{on\_analysis\_end()}.
  \item Add to a \texttt{DynamicObserverList} via \texttt{list.add<MyObs>(args...)}.
\end{enumerate}

\begin{verbatim}
template <typename ModelT>
class DriftObserver : public fall_n::ObserverBase<ModelT> {
    double max_drift_{0};
public:
    void on_step(const fall_n::StepEvent& ev,
                 const ModelT& model) override
    {
        double drift = compute_interstory_drift(ev.displacement, model);
        max_drift_ = std::max(max_drift_, drift);
    }
    double max_drift() const { return max_drift_; }
};
\end{verbatim}

\paragraph{Compile-time observer (for maximum performance):}

Implement the three methods (\texttt{on\_analysis\_start},
\texttt{on\_step}, \texttt{on\_analysis\_end}) as regular member functions
(no inheritance needed), then pass to \texttt{make\_composite\_observer}.


% ─────────────────────────────────────────────────────────────────────────
\subsection{Adding a New Incremental Control Policy}
\label{sec:arch:ext:control}

Any type satisfying \texttt{IncrementalControlPolicy} can be passed
to \texttt{solve\_incremental()} or \texttt{begin\_incremental()}:

\begin{verbatim}
struct MyPolicy {
    template <typename ModelT>
    void apply(double p, Vec f_full, Vec f_ext, ModelT* model) {
        // p ∈ [0,1] — must be absolute (idempotent for any p)
        VecCopy(f_full, f_ext);
        VecScale(f_ext, std::sin(p * M_PI / 2.0));  // sine ramp
    }
};

nl.solve_incremental(20, 4, MyPolicy{});
\end{verbatim}

\noindent
Alternatively, use \texttt{make\_control(lambda)} for one-off policies:

\begin{verbatim}
auto scheme = make_control([](double p, Vec f_full, Vec f_ext, auto* m) {
    VecCopy(f_full, f_ext);
    VecScale(f_ext, p * p);  // quadratic ramp
});
nl.solve_incremental(10, 4, scheme);
\end{verbatim}


% ─────────────────────────────────────────────────────────────────────────
\subsection{Adding a New StepDirector}
\label{sec:arch:ext:director}

Step directors are plain callables.  Use \texttt{pause\_on} or
\texttt{stop\_on} for quick construction, or write a full
\texttt{std::function}:

\begin{verbatim}
// Stop when inter-story drift exceeds 2%
auto dir = fall_n::step_director::stop_on(
    [](const fall_n::StepEvent& ev, const auto& model) {
        return max_interstory_drift(model) > 0.02;
    }
);
solver.set_director(dir);
\end{verbatim}


% ═══════════════════════════════════════════════════════════════════════════
% ═══════════════════════════════════════════════════════════════════════════
\section{Usage Tutorials}
\label{sec:arch:tutorials}


% ─────────────────────────────────────────────────────────────────────────
\subsection{Tutorial 1: Simple Linear Static Analysis}
\label{sec:arch:tut:linear}

\begin{verbatim}
#include "header_files.hh"

int main(int argc, char** argv) {
    PetscInitialize(&argc, &argv, nullptr, nullptr);
    {
        // 1. Build domain from Gmsh
        auto [domain, groups] = GmshDomainBuilder<3>::from_file("mesh.msh");

        // 2. Create model with elastic material
        auto material = make_elastic_3d(E, nu);
        Model<ThreeDimensionalMaterial> model(&domain);
        model.add_elements(material, groups["solid"]);
        model.apply_forces(groups["load"], {0, 0, -1000});
        model.constrain_nodes(groups["base"], {0, 1, 2});
        model.setup();

        // 3. Solve
        LinearAnalysis la(&model);
        la.solve();

        // 4. Post-process
        VTKModelExporter exporter("result.vtu");
        exporter.write(model);
    }
    PetscFinalize();
}
\end{verbatim}


% ─────────────────────────────────────────────────────────────────────────
\subsection{Tutorial 2: Nonlinear Incremental with Displacement Control}
\label{sec:arch:tut:pushover}

\begin{verbatim}
// Pushover analysis: node 42, direction X, target displacement 0.05 m
NonlinearAnalysis nl(&model);
DisplacementControl dc{42, 0, 0.05};
nl.solve_incremental(50, /*max_bisections=*/4, dc);
\end{verbatim}


% ─────────────────────────────────────────────────────────────────────────
\subsection{Tutorial 3: Dynamic Seismic Analysis with Observers}
\label{sec:arch:tut:seismic}

\begin{verbatim}
DynamicAnalysis dyn(&model);
dyn.set_density(2400.0);  // kg/m^3
dyn.set_rayleigh_damping(omega1, omega2, 0.05);  // 5% at two frequencies
dyn.set_boundary_conditions(bcs);
dyn.set_initial_conditions(bcs);

// Compose observers
auto obs = make_composite_observer<ModelT>(
    ConsoleProgressObserver<ModelT>{100},
    VTKSnapshotObserver<ModelT>{"output/", 50, profile, thickness},
    NodeRecorder<ModelT>{{42, 99}, {0, 1, 2}, 1}
);
dyn.set_observer(obs);

// Solve
dyn.solve(30.0, 0.01);  // 30 seconds, dt = 0.01 s

// Access recorded data
auto& recorder = obs.get<2>();
auto history = recorder.get_history(42, 0);  // node 42, dof 0
\end{verbatim}


% ─────────────────────────────────────────────────────────────────────────
\subsection{Tutorial 4: Multi-Phase Analysis with AnalysisDirector}
\label{sec:arch:tut:multiphase}

\begin{verbatim}
AnalysisDirector director;

// Phase 1: Static gravity preload
director.add_phase("gravity", [&](const AnalysisState& prev) {
    NonlinearAnalysis nl(&model);
    nl.solve_incremental(10);
    return nl.get_analysis_state();
});

// Phase 2: Dynamic seismic excitation
director.add_phase("seismic", [&](const AnalysisState& prev) {
    DynamicAnalysis dyn(&model);
    dyn.set_initial_displacement(prev.displacement);
    dyn.set_boundary_conditions(bcs);
    dyn.solve(30.0, 0.01);
    return dyn.get_analysis_state();
});

auto report = director.run();
report.print_summary();
\end{verbatim}


% ─────────────────────────────────────────────────────────────────────────
\subsection{Tutorial 5: Single-Step Control with StepDirector}
\label{sec:arch:tut:stepping}

\begin{verbatim}
NonlinearAnalysis nl(&model);
nl.begin_incremental(100);  // 100 load steps

// Step to half load, pause every 10 steps
auto dir = fall_n::step_director::pause_every_n<ModelT>(10);

while (nl.current_time() < 1.0 - 1e-14) {
    auto verdict = nl.step_to(1.0, dir);

    if (verdict == fall_n::StepVerdict::Pause) {
        // Inspect model between steps
        auto u_max = get_max_displacement(model);
        std::cout << "p = " << nl.current_time()
                  << ", u_max = " << u_max << "\n";

        if (u_max > threshold) break;  // early termination
    }

    if (verdict == fall_n::StepVerdict::Stop) {
        std::cerr << "Analysis diverged!\n";
        break;
    }
}
\end{verbatim}


% ═══════════════════════════════════════════════════════════════════════════
% ═══════════════════════════════════════════════════════════════════════════
\section{Cross-Module Dependency Map}
\label{sec:arch:crossmodule}

The full framework follows a strict layered dependency:

\begin{center}
\begin{tikzpicture}[
    layer/.style={draw, rounded corners=3pt, minimum width=14cm,
                  minimum height=1.0cm, font=\small},
    arr/.style={-{Stealth[length=5pt]}, thick}
  ]
  \node[layer, fill=gray!5]   (L6) at (0, 7.0)
    {\textbf{Post-processing} : VTK exporters, state queries};
  \node[layer, fill=blue!8]   (L5) at (0, 5.5)
    {\textbf{Analysis} : LinearAnalysis, NLAnalysis, DynamicAnalysis, AnalysisDirector};
  \node[layer, fill=green!8]  (L4) at (0, 4.0)
    {\textbf{Model} : Model, BoundaryCondition, NodeSelector, GroundMotionRecord};
  \node[layer, fill=orange!8] (L3) at (0, 2.5)
    {\textbf{Domain + Elements} : Domain, ContinuumElement, BeamElement, ShellElement};
  \node[layer, fill=red!8]    (L2) at (0, 1.0)
    {\textbf{Materials + Continuum} : Material, ConstitutiveRelation, Tensor, KinematicPolicy};
  \node[layer, fill=purple!5] (L1) at (0,-0.5)
    {\textbf{Geometry + Numerics} : ElementGeometry, Quadrature, Shape functions};

  \draw[arr] (L1) -- (L2);
  \draw[arr] (L2) -- (L3);
  \draw[arr] (L3) -- (L4);
  \draw[arr] (L4) -- (L5);
  \draw[arr] (L5) -- (L6);
\end{tikzpicture}
\end{center}

No layer may depend on a higher layer.  The only allowed cross-layer
access is downward.  The \texttt{header\_files.hh} aggregator includes
all layers in bottom-up order.


% ═══════════════════════════════════════════════════════════════════════════
\section{Ontological Critique}
\label{sec:arch:ontology}

This section critically examines the semantic coherence of the
framework's naming conventions.

\paragraph{Well-aligned names.}
\begin{itemize}[nosep]
  \item \texttt{NonlinearAnalysis} / \texttt{DynamicAnalysis} /
        \texttt{LinearAnalysis}: clear solver-type distinction.
  \item \texttt{SteppableSolver}: accurately describes the stepping
        protocol (not ``Solver'' alone, which is too generic).
  \item \texttt{AnalysisDirector}: orchestrator pattern, correctly
        signalling ``directs'' rather than ``is.''
  \item \texttt{StepDirector}: breakpoint evaluator, distinct from
        \texttt{AnalysisDirector} in scope.
  \item \texttt{IncrementalControlPolicy}: ``control'' + ``policy'' maps
        precisely to the Strategy pattern.
  \item \texttt{NodeSelector}: composable predicate --- ``selector'' is a
        well-established pattern name.
  \item \texttt{BoundaryConditionSet}: ``set'' correctly denotes a
        collection of BCs, not a single BC.
\end{itemize}

\paragraph{Areas for future refinement.}
\begin{itemize}[nosep]
  \item \texttt{DoF.hh} / \texttt{DoF\_Interface} / \texttt{DoF\_Handler}: the
        ``Interface'' / ``Handler'' distinction is unclear.  Superseded by
        \texttt{DoFStorage} (clean policy-based design); remove when possible.
  \item \texttt{ModelBuilder}: dormant.  The name conflicts with the active
        \texttt{StructuralModelBuilder} and \texttt{BuildingDomainBuilder}.
        Remove or rename to \texttt{LegacyModelBuilder} before reuse.
  \item \texttt{Observers.hh}: the filename is generic; consider
        \texttt{ConcreteObservers.hh} or splitting into per-concern files
        if the file grows further.
  \item \texttt{ElemPolicy} template parameter: ``Policy'' is accurate for
        the storage strategy, but could be confused with
        \texttt{MaterialPolicy} or \texttt{KinematicPolicy} which govern
        physics.  A name like \texttt{StoragePolicy} would reduce ambiguity.
\end{itemize}


% ═══════════════════════════════════════════════════════════════════════════
\section{Code Cleanup Summary}
\label{sec:arch:cleanup}

The following changes were applied as part of this architectural review:

\begin{enumerate}
  \item \textbf{\texttt{Analysis.hh}}: Removed stale \texttt{ANALYSIS CONCEPT : TODO}
        comment.  Reordered includes into five labelled groups following
        the dependency graph: concepts/types $\to$ damping $\to$
        solvers $\to$ orchestration $\to$ observers.  Added module-level
        documentation comment.

  \item \textbf{\texttt{DoF.hh}}: Fixed include-guard typo
        (\texttt{DEGGREE} $\to$ \texttt{DEGREE}).

  \item \textbf{\texttt{header\_files.hh}}: Updated the \texttt{DoF.hh}
        transition comment to match current reality.

  \item \textbf{\texttt{ModelBuilder.hh}}: Added a \texttt{DORMANT} banner
        explaining the class's status, its historical role, and the active
        replacements (\texttt{StructuralModelBuilder},
        \texttt{GmshDomainBuilder}).
\end{enumerate}

All 52 tests pass after these changes.


% ═══════════════════════════════════════════════════════════════════════════
\section{Summary}
\label{sec:arch:summary}

The Analysis and Model modules form a coherent, extensible architecture:

\begin{itemize}[nosep]
  \item \textbf{Analysis}: concepts define contracts
        (\texttt{SteppableSolver}, \texttt{IncrementalControlPolicy});
        template engines wrap PETSc solvers with RAII; observers decouple
        recording; \texttt{AnalysisDirector} composes multi-phase workflows.

  \item \textbf{Model}: the \texttt{Model<>} template aggregates domain,
        elements, and PETSc structures; \texttt{BoundaryConditionSet}
        provides a declarative BC API; \texttt{NodeSelector} enables
        composable queries; builders specialise construction for
        structural and building models.

  \item \textbf{Extension}: new solvers, observers, control policies, and
        step directors can be added without modifying existing code ---
        the Open/Closed principle is consistently respected.

  \item \textbf{Ontology}: naming is generally precise and well-aligned;
        legacy artifacts (\texttt{DoF.hh}, \texttt{ModelBuilder.hh}) are
        documented and quarantined for future removal.
\end{itemize}
